\section{指标体系：真人连接率之外还要看什么}

Tristan 提到北极星是真人连接率，这是非常准确的方向。但一个社交产品不能只看单一指标。真人连接率要配合留存、误代表率、context 完成度、隐私撤回率、AI/真人互动比例和长期关系质量一起看。

\subsection{指标表}

\noindent\begin{longtable}{p{0.24\textwidth}p{0.32\textwidth}p{0.35\textwidth}}
\toprule
指标 & 衡量什么 & 误用风险 \\
\midrule
真人连接率 & AI 促成的真人有效互动比例 & 可能诱导过度打扰。 \\
真人行为比例 & 网络中真人而非 AI 的互动占比 & 过低会变成 AI 自嗨。 \\
Context 完成度 & 用户是否提供足够可用 context & 过度追求会侵犯隐私。 \\
误代表率 & 分身输出被用户否定的比例 & 需要精细记录和快速修正。 \\
留存 & 用户是否持续回来 & 可能被虚假热闹抬高。 \\
关系质量 & 连接是否持续、有价值、低伤害 & 难量化但最接近愿景。 \\
\bottomrule
\end{longtable}

\begin{knowledgebox}{为什么指标需要成组看}
只看真人连接率，可能鼓励粗暴拉人；只看留存，可能鼓励上瘾；只看 AI 互动数，可能制造虚假繁荣。AI 社交需要一组互相制衡的指标。
\end{knowledgebox}

\subsection{本章小结}

真人连接率是北极星，但不是唯一仪表盘。好的 AI 社交指标体系要同时看连接、隐私、误代表和关系质量。

